% 数学模式 CJK 兜底：\text{中文} 走文本模式交给 xeCJK（unicode-math 的 \text 缺 CJK 字形）
\usepackage{unicode-math}
\setmathfont{latinmodern-math.otf}
% 页边距：左右 2.2cm、上下 2.4cm（正文可用宽度 ~165mm，缓解窄列表格溢出）
\usepackage[top=2.4cm,bottom=2.4cm,left=2.2cm,right=2.2cm]{geometry}
% 表格排版：字号比正文降一档（\small）、行末放宽应急间距、列间距收紧、长标识符允许在下划线后断行
\usepackage{etoolbox}
\usepackage{ragged2e}
\AtBeginEnvironment{longtable}{\small\RaggedRight\setlength{\emergencystretch}{2em}}
\AtBeginEnvironment{tabular}{\small\RaggedRight}
\renewcommand{\tabcolsep}{4pt}
\makeatletter
\renewcommand{\_}{\discretionary{}{}{\textunderscore}}
\makeatother
% \text 重定义：数学模式内的中文文本交给 xeCJK 处理
\makeatletter
\renewcommand{\text}[1]{\textnormal{#1}}
\makeatother
% 超宽浮动限制
\usepackage{graphicx}
\PassOptionsToPackage{hyphens}{url}
% 放宽浮动体高度限制（消除 Float too large：默认页顶浮动限 0.7 页高，专著含大量纵向示意图）
\renewcommand{\topfraction}{0.93}
\renewcommand{\bottomfraction}{0.93}
\renewcommand{\textfraction}{0.04}
\setcounter{topnumber}{4}
\setcounter{bottomnumber}{4}
\setcounter{totalnumber}{6}
\setcounter{dbltopnumber}{4}

% 章级标题（\section，对应 pandoc 一级标题）强制另起一页
\let\oldsection\section
\renewcommand{\section}{\clearpage\oldsection}
